% ---------------------------------------------------------------------------
% Kapitel 5: Implementierung des Starship-Gehwegroboters
% Ersetzt den Platzhaltertext von Kapitel 5. Aufbau wie Kapitel 6 (Clevon).
% Abbildungen: figures/starship/*  |  Literatur: starship_quellen.bib
% Stand: 22.10.2026, Repository-Branch merge/clevon-starship-ein-repo
% ---------------------------------------------------------------------------
\chapter{Implementierung des Starship-Gehwegroboters}
\label{chap:starship}

Dieses Kapitel beschreibt die Modellierung und Integration des Starship-Gehwegroboters
in CARLA. Im Mittelpunkt stehen die Fahrzeuggeometrie, die Umsetzung der Fahrphysik mit
sechs simulierten Rädern unter Chaos Vehicles, die Sensorik sowie ein Bordstein-Assistent,
der das Überfahren von Bordsteinen ermöglicht. Beschrieben wird der abschließende
Implementierungsstand unter CARLA 0.10.0 und Unreal Engine 5.5.4; die ROS-2-Anbindung
folgt in Kapitel~7.

\section{Referenzfahrzeug und Modellumfang}
\label{sec:starship_referenz}

Der Starship-Roboter der Starship Technologies OÜ bildet den gehwegbasierten
Lieferroboter des Projekts ab. Er fährt mit höchstens 6\,km/h, wiegt leer etwa 25\,kg
und nimmt bis zu 9,1\,kg Zuladung auf \parencite{wikipedia_starship}. Das Fahrwerk besteht
aus sechs Rädern. Das Fahrwerkspatent des Herstellers nennt mindestens vier Elektromotoren,
etwa je einen für jedes Vorderrad und je einen für Mittel- und Hinterrad einer Seite; die
Motoren können auch in den Rädern sitzen \parencite{starship_patent_fahrwerk}. Ein weiteres Patent beschreibt Räder mit etwa
20\,cm Durchmesser, die an Kipphebeln sitzen und Hindernisse von 15 bis 20\,cm Höhe
überwinden \parencite{starship_patent_bogie}. Ein Herstellerdatenblatt mit Abmessungen,
Motorleistung oder Sensortypen ist nicht veröffentlicht. Die Modellierung stützt sich
deshalb auf diese Angaben, auf Patente des Herstellers und auf Presseberichte.

Das Simulationsmodell umfasst die Fahrzeuggeometrie mit Deckel und Antenne,
Fahr- und Lenkfunktionen mit sechs angetriebenen Rädern, einen Bordstein-Assistenten
sowie 26 virtuelle Sensoren. Lichtfunktionen und ein öffnender Deckel sind nicht
Bestandteil der Umsetzung.

\subsection{Eigene Modellierung und HKA-Design}
\label{sec:starship_design}

Der Roboter wurde in Blender eigenständig nachmodelliert. Grundform und Proportionen,
also kastenförmiger Rumpf mit abgerundeten Kanten, Deckel, sechs Räder in drei Achsen
und eine Signalantenne, folgen dem Original. Um keine geschützte Gestaltung zu
übernehmen, wurde das Erscheinungsbild bewusst verändert: Der Rumpf ist anthrazit, auf
beiden Seiten steht der Schriftzug „HKA“ in Rot, die Front trägt ein rotes Lichtband und
die Antenne endet in einem roten Signallicht statt einer Fahne
(Abbildung~\ref{fig:starship_hka}). Ein Hersteller-CAD-Modell wurde nicht verwendet; alle
Bauteile sind eigene Modellierung.

\begin{figure}[htbp]
  \centering
  \includegraphics[width=0.62\textwidth]{figures/starship/hka_left.png}
  \caption{Eigenes Starship-Modell im HKA-Design. Eigene Darstellung aus Blender.}
  \label{fig:starship_hka}
\end{figure}

Sichtbare Sensorgehäuse sind im Modell nicht ausgearbeitet. Die Messfunktion entsteht
ausschließlich durch die zur Laufzeit angebrachten virtuellen Sensoren in CARLA
(Abschnitt~\ref{sec:starship_sensorik}).

\section{Geometrie, Rig und CARLA-Integration}
\label{sec:starship_geometrie}

Das Fahrzeugskelett besteht aus dem Basisknochen \texttt{VehicleBase} und sechs
Radknochen \texttt{Wheel\_\{Front,Middle,Rear\}\_\{Left,Right\}}, die jeweils im
Radmittelpunkt liegen (Abbildung~\ref{fig:starship_rig}). Rumpf und Räder sind starr an
ihre Knochen gebunden. Alle Knochen zeigen in Blender in $+y$-Richtung und haben keinen
Rollwinkel. Damit stimmen ihre lokalen Achsen in Unreal Engine mit den Fahrzeugachsen
überein. Ein um 90° gekippter Basisknochen hatte im ersten Stand dazu geführt, dass die
Fahrwerksberechnung die Räder seitlich abtastete und kein Rad Bodenkontakt meldete.

\begin{figure}[htbp]
  \centering
  \includegraphics[width=0.55\textwidth]{figures/starship/blender_bones.png}
  \caption{Basisknochen \texttt{VehicleBase} und sechs Radknochen im Starship-Modell
  (Vorstufe mit gleichem Skelett, noch mit Fahne). Eigene Darstellung aus Blender.}
  \label{fig:starship_rig}
\end{figure}

Die Modellquellen liegen im Repository unter \texttt{02\_starship\_emodell/models}. Die
Datei \texttt{starship\_hka\_v5.blend} enthält den aktuellen Stand, der Export
\texttt{starship\_hka\_v5.fbx} ist die Grundlage für den Import. Der FBX-Export verwendet
Zentimeter; Mesh- und Knochenskalierung betragen jeweils 1/1/1, der Importfaktor ist 1.
Ein früherer Export mit Knotenskalierung 100 ließ den Roboter in CARLA durch die Fahrbahn
fallen, ein Export in Metern kam hundertfach verkleinert an.

Die Unreal-Engine-4-Assets des Ausgangsstands ließen sich in Unreal Engine 5 nicht
öffnen. Das Fahrzeug wurde deshalb aus der FBX-Datei und den ausgelesenen Fahrzeugwerten
des Ausgangsstands neu aufgebaut. Alle Schritte laufen als Python-Skripte im geöffneten
Editor (Remote Execution) und prüfen ihr Ergebnis selbst: Mesh-Import, Physics Asset,
Animations-Blueprint, sechs Rad-Blueprints, Fahrzeug-Blueprint und Registrierung
\parencite{rauch_lee_repo}. Der Aufbau ist damit vollständig reproduzierbar; fertige
Unreal-Engine-5-Assets liegen nicht im Repository.

Der Fahrzeug-Blueprint \texttt{BP\_Starship} basiert auf \texttt{BaseVehiclePawnNW},
der CARLA-Basisklasse für Fahrzeuge mit mehr als vier Rädern, die auch der
sechsrädrige Lkw \texttt{EuropeanHGV} verwendet. Seine Registrierung erfolgt über die
\texttt{VehicleParameters.json} unter der Kennung \texttt{vehicle.starship.bot} mit sechs
Rädern. Die Kennung ist dieselbe wie im Ausgangsstand, sodass bestehende Clients
unverändert weiterlaufen.

\subsection{Abmessungen}
\label{sec:starship_masse}

Tabelle~\ref{tab:starship_masse} vergleicht die verfügbaren Referenzangaben mit dem
Modell. Die Modellmaße stammen aus der von CARLA gemeldeten Actor-Bounding-Box und den
Radknochen. Da der Hersteller keine Außenmaße veröffentlicht, ist nur ein grober
Vergleich mit der im Kamerapatent genannten Robotergröße möglich.

\begin{table}[htbp]
  \centering
  \caption{Referenzangaben und Modellwerte des Starship. Referenz: \textcite{wikipedia_starship},
  \textcite{starship_patent_kamera}, \textcite{starship_patent_bogie}; Modell: eigene Messung.}
  \label{tab:starship_masse}
  \begin{tabular}{llll}
    \hline
    Größe & Referenz & Modell & Bemerkung \\
    \hline
    Länge $\times$ Breite & etwa 60\,cm Robotergröße & 0,564 $\times$ 0,554\,m & Bounding-Box \\
    Höhe mit Antenne & -- & 0,658\,m & Bounding-Box \\
    Rumpf (Kollisionsbox) & -- & 0,55 $\times$ 0,45 $\times$ 0,30\,m & 19 bis 49\,cm über Boden \\
    Raddurchmesser & etwa 20\,cm & 18\,cm & Radius 9\,cm \\
    Radstand / Spurweite & -- & 0,382 / 0,512\,m & aus Radknochen \\
    Anzahl Räder & 6 & 6 & alle simuliert \\
    Masse & 25\,kg leer, bis 9,1\,kg Zuladung & 35\,kg & beladen \\
    Höchstgeschwindigkeit & 6\,km/h & 6,3\,km/h & gemessen \\
    \hline
  \end{tabular}
\end{table}

\subsection{Aufgaben der Blueprint-Bausteine}
\label{sec:starship_bausteine}

Das Fahrzeug entsteht wie der Clevon aus mehreren Assets. Das Skeletal Mesh
\texttt{SK\_Starship} enthält Geometrie und Skelett, das Physics Asset den
Kollisionskörper. Sechs Chaos-Wheel-Blueprints \texttt{BP\_Starship\_\{FLW, FRW, MLW,
MRW, RLW, RRW\}} legen die radbezogenen Parameter fest. Die Chaos-Fahrzeugbewegungs\-komponente
berechnet die Fahrzeugbewegung, der Animations-Blueprint \texttt{AnimBP\_Starship}
überträgt die Radbewegungen auf das Skelett. \texttt{BP\_Starship} verbindet diese
Bausteine und enthält zusätzlich die Komponente \texttt{CurbAssist}
(Abschnitt~\ref{sec:starship_bordstein}). Eine eigene Blueprint-Logik im Event Graph ist
nicht erforderlich, da weder Lichter noch bewegliche Zusatzteile angesteuert werden.

Zwei Einstellungen weichen von der Basisklasse ab und sind für die Funktion notwendig.
\texttt{BaseVehiclePawnNW} liefert den Fahrzeugkörper ohne Physiksimulation aus;
ohne Aktivierung von \emph{Simulate Physics} bleibt der Roboter an der Spawnposition
eingefroren. Außerdem muss das erste Wheel Setup ein lenkendes Rad sein, da CARLA den
maximalen Lenkwinkel aus Rad~0 liest. Die von CARLA genutzte Box
\texttt{VehicleBounds} passt CARLA beim Spawnen selbst an das Mesh an.

% TODO: Screenshot der Komponentenstruktur von BP_Starship in Unreal Engine 5.5.4
% (analog Abbildung 6.4) am Linux-Rechner aufnehmen und hier einfügen.

\subsection{Skelettanimation und Kollisionskörper}
\label{sec:starship_animation}

Der Animations-Blueprint ist eine Kopie des Blueprints \texttt{AnimBP\_EuropeanHGV}, auf
das Starship-Skelett umgestellt. Der Animationsgraph verbindet Mesh Space Ref Pose,
Wheel Controller, Component To Local und Output Pose. Der Wheel Controller erhält
Radknochen, Raddrehung, Lenkwinkel und Federweg zur Laufzeit aus den Wheel Setups der
Bewegungskomponente und ist daher nicht an die Knochennamen des Lkw gebunden. Im
Ausgangsstand unter PhysX waren nur vier Räder simuliert; die Mittelräder drehten über
zwei Copy-Bone-Knoten lediglich optisch mit. Unter Chaos werden alle sechs Räder
simuliert, die Copy-Bone-Knoten entfallen.

Das Physics Asset enthält genau einen Körper: eine Box auf \texttt{VehicleBase} mit
55\,$\times$\,45\,$\times$\,30\,cm, deren Unterkante 19\,cm und Oberkante 49\,cm über dem
Boden liegen. Radkörper gibt es nicht. Die Rad-Boden-Kontakte ermittelt die
Chaos-Fahrzeugbewegungskomponente wie beim Clevon mit einer Spherecast-Abfrage je Rad, die
den Radradius berücksichtigt. Zunächst war ein senkrechter Strahl (Raycast) eingestellt,
weil eine Kugelabfrage die Bordsteinkante seitlich trifft und die Federkraft dann entlang
der nahezu waagerechten Flächennormale nach hinten wirkt. Im Fahrbetrieb mit
Bordstein-Assistent zeigte sich jedoch, dass das Rad mit Raycast in die Bordsteinfront
eintaucht und dort einhakt, da der Strahl nur einen Punkt unter der Radmitte prüft. Mit
Spherecast erfasst die Abfrage die Kante bereits mit der Radkontur, und das Rad steigt
auf. Die Radabfrage ist deshalb auf Spherecast umgestellt.
Die Unterkante der Box wurde von ursprünglich 16 auf 19\,cm angehoben,
damit der Rumpf 15\,cm hohe Bordsteine auch dann nicht berührt, wenn die Vorderräder
bereits oben sind.

\subsection{Lenkung}
\label{sec:starship_lenkung}

Der reale Starship lenkt nach seinem Fahrwerkspatent wie ein Kettenfahrzeug: Die
Drehachsen der Räder bleiben auch bei Kurvenfahrt fest zum Roboter, die Seiten werden
unterschiedlich schnell angetrieben \parencite{starship_patent_fahrwerk}. Das Modell
verwendet stattdessen eine Achsschenkellenkung mit Ackermann-Geometrie: Die Vorderräder
schlagen bis 30° ein, Mittel- und Hinterräder sind nicht lenkbar. Gründe sind die
Fahrzeugschnittstelle von CARLA, deren Steuerbefehl einen Lenkwert statt zweier
Seitengeschwindigkeiten vorsieht, sowie die Fahrwerksmodelle von PhysX und Chaos, die
gelenkte Räder voraussetzen. Fahrclient, Autopilot und ROS-2-Steuerbefehl
(Kapitel~7) bleiben dadurch ohne Anpassung nutzbar. Ein Wenden auf der Stelle ist
damit nicht möglich; der gemessene Wenderadius beträgt 2,2 bis 2,4\,m.

\section{Fahrphysik und Antrieb}
\label{sec:starship_fahrphysik}

Die Fahrphysik wird durch Chaos Vehicles berechnet. Die Parametrierung umfasst Masse und
Schwerpunkt, Rad- und Federungseigenschaften sowie das Antriebsmodell. Die gewählten
Werte beschreiben einen Simulationsbetriebsfall und keine Herstellerkalibrierung.

\subsection{Masse, Schwerpunkt und Fahrwerk}
\label{sec:starship_masse_fahrwerk}

Die Simulation verwendet eine Fahrzeugmasse von 35\,kg. Sie entspricht der
Leermasse von etwa 25\,kg zuzüglich der maximalen Zuladung von 9,1\,kg
\parencite{wikipedia_starship} und beschreibt damit den beladenen Fahrzeugzustand.
Am Fahrzeug-Mesh ist ein Schwerpunktversatz von $(0, 0, -15)$\,cm eingestellt. Ein aus
einem Pkw-Blueprint übernommener Versatz von $(25, 0, -30)$\,cm hatte den Schwerpunkt vor
die Vorderachse gelegt, die nur 19\,cm vor der Fahrzeugmitte liegt.

\begin{table}[htbp]
  \centering
  \caption{Wesentliche Geometrie- und Fahrwerksparameter des Starship-Modells unter
  Chaos Vehicles. Die Einstellungen stellen keine Herstellerkalibrierung dar.}
  \label{tab:starship_fahrwerk}
  \begin{tabular}{ll}
    \hline
    Parameter & Wert \\
    \hline
    Fahrzeugmasse & 35\,kg \\
    Schwerpunktversatz & $(0, 0, -15)$\,cm \\
    Radstand / Spurweite & 0,382 / 0,512\,m \\
    Radradius / Radbreite / Radmasse & 0,09\,m / 0,04\,m / 1\,kg \\
    Maximaler Lenkwinkel vorne / Mitte / hinten & 30° / 0° / 0° \\
    Lenkgeometrie & Ackermann \\
    Angetriebene Räder & alle sechs \\
    Maximales Bremsmoment je Rad & 20\,Nm \\
    Feststellbremsmoment Mittel- und Hinterräder & 40\,Nm \\
    Maximaler Einfederweg / Ausfederweg & 18 / 4\,cm \\
    Federrate & 14,3\,N/cm \\
    Suspension Damping Ratio & 0,7 \\
    Friction Force Multiplier & 2,0 \\
    Rad-Boden-Abfrage & Spherecast \\
    \hline
  \end{tabular}
\end{table}

Die Vorderachse übernimmt die Lenkung. Die Betriebsbremse wirkt auf alle sechs Räder, die
Feststellbremse auf Mittel- und Hinterräder. Der Radradius von 9\,cm wurde am Modell
gemessen; Radmasse, Bremsmomente und Federwege sind Annahmen für einen kleinen Roboter.

Die Federung ist auf Bordsteine ausgelegt. Der große Einfederweg von 18\,cm erlaubt
einem Rad, auf eine 15\,cm hohe Kante zu steigen, ohne auf den harten Anschlag zu treffen.
Die Federrate ergibt sich aus der Radlast und dem Ausfederweg. Bei 35\,kg auf sechs
Rädern beträgt die Radlast $F_\mathrm{R} = 35\,\mathrm{kg} \cdot 9{,}81\,\mathrm{m/s^2} / 6
\approx 57\,\mathrm{N}$, die Federrate
\begin{equation}
  k = \frac{F_\mathrm{R}}{s_\mathrm{aus}} = \frac{57\,\mathrm{N}}{4\,\mathrm{cm}}
  \approx 14{,}3\,\mathrm{N/cm}.
  \label{eq:starship_feder}
\end{equation}
Damit federt jedes Rad in Ruhe genau um den Ausfederweg ein, und die Radknochen stehen
auf der modellierten Höhe. Die Vorspannung wird dafür nicht benötigt, da Chaos sie in der
verwendeten Version nicht auswertet. Die Eigenfrequenz liegt bei etwa 2,5\,Hz; das macht
das Fahrwerk auf Fahrbahnmarkierungen und kleinen Unebenheiten weich. Chaos rechnet die
eingegebene Federrate intern mit dem Faktor 100 um, der Eingabewert entspricht deshalb
der Federsteifigkeit in N/cm. Federrate, Dämpfung und Reifenkräfte sind nicht anhand
von Messungen am realen Fahrzeug kalibriert.

\subsection{Antriebsmodell mit sechs angetriebenen Rädern}
\label{sec:starship_antrieb}

Beim realen Starship treiben mehrere Elektromotoren alle sechs Räder an, die Motoren
können in den Rädern sitzen \parencite{starship_patent_fahrwerk}. Chaos bildet dagegen
einen zentralen Motor mit Getriebe und Differenzial ab. Das Modell
verwendet deshalb einen Ersatzmotor mit flacher Drehmomentkennlinie, dessen Moment über
ein Allraddifferenzial auf alle sechs Räder verteilt wird. Die Verteilung erfolgt nach
Achstyp: Die beiden Vorderräder erhalten je 25\,\%, die vier Mittel- und Hinterräder je
12,5\,\%. Leistung und Drehmoment des realen Antriebs sind nicht veröffentlicht.

Die Höchstgeschwindigkeit wird über die maximale Motordrehzahl begrenzt, nicht über das
Drehmoment. Mit $n_\mathrm{max} = 720\,\mathrm{min^{-1}}$, der Gangübersetzung
$i_\mathrm{G} = 1$, der Achsübersetzung $i_\mathrm{A} = 4$ und dem Radradius
$r = 0{,}09\,\mathrm{m}$ ergibt sich
\begin{equation}
  v_\mathrm{max} = \frac{n_\mathrm{max}}{i_\mathrm{G}\, i_\mathrm{A}} \cdot \frac{2\pi r}{60}
  = \frac{720}{4} \cdot \frac{2\pi \cdot 0{,}09\,\mathrm{m}}{60\,\mathrm{s}}
  \approx 1{,}70\,\mathrm{m/s} \approx 6{,}1\,\mathrm{km/h}.
  \label{eq:starship_vmax}
\end{equation}
Dabei ist $n_\mathrm{max}$ in $\mathrm{min^{-1}}$ einzusetzen. Gemessen wurden 6,3\,km/h.
Das maximale Motormoment beträgt 3\,Nm. Mit einem angenommenen Wirkungsgrad von 0,9
ergibt sich eine Zugkraft von
\begin{equation}
  F_\mathrm{Z} = \frac{M_\mathrm{max}\, i_\mathrm{G}\, i_\mathrm{A}\, \eta}{r}
  = \frac{3\,\mathrm{Nm} \cdot 4 \cdot 0{,}9}{0{,}09\,\mathrm{m}} = 120\,\mathrm{N},
  \label{eq:starship_zugkraft}
\end{equation}
rechnerisch also bis zu 3,4\,m/s² Beschleunigung bei 35\,kg. In Chaos wird die
Kennlinie durch das Kurven-Asset \texttt{Curve\_Starship\_Torque} mit dem Wert 1,0 über
den gesamten Drehzahlbereich und den Parameter Max Torque abgebildet. Ohne
Drehmomentkurve schaltet Chaos die Antriebssimulation ab.

Im Ausgangsstand unter PhysX trieb ein Vorderradantrieb mit 10 bis 15\,Nm und einer
Achsübersetzung von 3,4 nur zwei Räder an. Der Antrieb aller sechs Räder entspricht dem
Original und verbessert die Traktion an Kanten; Motormoment und Übersetzung wurden so
angepasst, dass die Höchstgeschwindigkeit weiterhin bei rund 6\,km/h liegt.

\section{Sensorik und Parametrierung}
\label{sec:starship_sensorik}

Das Sensorkonzept umfasst 26 Sensoren: neun RGB-Kameras, drei ToF-Tiefenkameras, vier
Radare, acht Ultraschallsensoren, GNSS und IMU. Anzahl von Kameras, Radaren und
Ultraschallsensoren entspricht den Angaben zum Starship in
\textcite{mit_technology_review_2018}. Die Anordnung der Kameras folgt dem Kamerapatent
des Herstellers \parencite{starship_patent_kamera}. Für ToF, Radar, Ultraschall, GNSS und
IMU sind keine Einbaupositionen veröffentlicht; ihre Lage ist aus der Aufgabe des
jeweiligen Sensors abgeleitet. Neuere Angaben des Herstellers nennen zwölf Kameras
einschließlich der ToF-Kameras \parencite{starship_robots}, was mit neun RGB- und drei
ToF-Kameras übereinstimmt.

Die virtuellen Sensoren werden als CARLA-Actors an das Fahrzeug gebunden. Positionen,
Orientierungen und Attribute sind je Sensortyp in JSON-Profilen hinterlegt, die sowohl
die Anzeigeprogramme als auch die ROS-Sensorskripte verwenden \parencite{rauch_lee_repo}.
Abbildung~\ref{fig:starship_sensoren} zeigt die Anordnung, Tabelle~\ref{tab:starship_sensoren}
fasst die Profile zusammen. Die Messraten sind Sollwerte in Simulationszeit; die
Referenzprodukte dienen der Parametrierung und belegen nicht die Identität der im realen
Starship verbauten Sensoren.

\begin{figure}[htbp]
  \centering
  \includegraphics[width=\textwidth]{figures/starship/starship_sensoranordnung.png}
  \caption{Anordnung der 26 Sensoren aus den gespeicherten CARLA-Montagewerten.
  K: Kamera, T: ToF, R: Radar, U: Ultraschall, G: GNSS, I: IMU. GNSS und IMU liegen bei
  $x = y = 0$ und sind in der Draufsicht nur zur Lesbarkeit versetzt. Pfeillängen
  schematisch. Eigene Darstellung.}
  \label{fig:starship_sensoren}
\end{figure}

\begin{table}[htbp]
  \centering
  \caption{Sensorprofile des Starship-Modells. Messraten sind Sollwerte in Simulationszeit.}
  \label{tab:starship_sensoren}
  \begin{tabular}{p{0.2\textwidth}p{0.42\textwidth}p{0.3\textwidth}}
    \hline
    Sensor & Wesentliche Parameter & Referenz \\
    \hline
    RGB, 9 Stück & 640 $\times$ 480 Pixel (seitlich hochkant 480 $\times$ 640); 80° bzw. 60° horizontal; 4\,Hz
      & Kamerapatent Starship \\
    ToF, 3 Stück & Tiefenkamera 320 $\times$ 240 Pixel; 90°; 5\,m; 10\,Hz & ESPROS epc660 \\
    Radar, 4 Stück & 20\,m; 60° $\times$ 20°; 1500 Strahlen/s; 20\,Hz & TI IWR6843AOP \\
    Ultraschall, 8 Stück & Hinderniserkennung 3\,m; Radius 0,15\,m; 20\,Hz & MaxBotix MB1242 \\
    GNSS & 5\,Hz; Streuung $\sigma \approx 2$\,m horizontal, 1\,m Höhe & Empfänger ohne RTK (Annahme) \\
    IMU & 100\,Hz; $\sigma$ 0,02\,m/s² bzw. 0,002\,rad/s & MEMS-IMU (Annahme) \\
    \hline
  \end{tabular}
\end{table}

Zur Verringerung der Renderlast erzeugen Kamera- und ToF-Anzeige jeweils nur die Sensoren
des ausgewählten Reiters; alle Positionen bleiben im Profil hinterlegt. Mit allen neun
Kameras und drei ToF-Kameras gleichzeitig stürzte der Editor bei 16\,GB Grafikspeicher
ab. Die Anzeigeprogramme für Kameras, ToF, Radar und Ultraschall sowie die GNSS- und
IMU-Erweiterung des Fahrclients wurden unter CARLA 0.10.0 funktional erprobt.

\subsection{Kameras und Blickrichtungen}
\label{sec:starship_kameras}

Das Kamerapatent beschreibt eine Einzelkamera vorne mit waagerechter Sicht sowie
Stereopaare vorne, links, rechts und hinten \parencite{starship_patent_kamera}. Die
Einzelkamera sitzt vorne mittig auf 42\,cm Höhe. Das vordere Paar liegt darunter auf
40\,cm, hat eine Stereobasis von 8\,cm und ist um 30° nach unten geneigt; das Patent nennt
5 bis 10\,cm Basis und 20 bis 40° Neigung, damit Bordstein und Boden direkt vor dem
Roboter erfasst werden. Die seitlichen Paare haben eine Basis von 18\,cm (Patent: 15 bis
20\,cm), das hintere Paar 8\,cm. Seitliche und hintere Kameras sind um 10° nach unten
geneigt. Die Kamerahöhe ist im Patent nicht angegeben und wurde knapp unter die
Deckelkante gelegt, wo die Sicht am weitesten reicht und die Räder sie nicht verdecken.

Das Patent nennt ein Sichtfeld von etwa 60° $\times$ 80°, hochkant eingebaute
Seitenkameras und eine Bildrate von etwa vier Bildern pro Sekunde. Das Modell übernimmt
dies mit 80° horizontal für Front- und Heckkameras, 60° horizontal für die hochkant
eingebauten Seitenkameras (480 $\times$ 640 Pixel) und einem Messintervall von 0,25\,s.
Die Auflösung ist im Patent nicht angegeben und als Annahme auf 640 $\times$ 480 Pixel
gesetzt. Linsenverzeichnung wird nicht nachgebildet.

\subsection{ToF-Tiefenkameras}
\label{sec:starship_tof}

Starship verwendet ToF-Bildsensoren des Typs ESPROS epc660 \parencite{robot_report_espros}.
Dieser Chip hat 320 $\times$ 240 Bildpunkte \parencite{espros_epc660}. Das Modell bildet
die ToF-Sensoren nach Abschnitt~3.2 durch CARLA-Tiefenkameras mit dieser Auflösung, 90°
Sichtfeld und einer auf 5\,m begrenzten Messdistanz nach. Drei Sensoren sitzen vorne und
an beiden Seiten auf 30\,cm Höhe und sind um 15° nach unten geneigt. Der Mittelstrahl
trifft den Boden damit etwa 1,1\,m vor dem Sensor; erfasst werden Bordsteine, Stufen und
Löcher im Nahbereich. Nach hinten decken Stereokameras, Radar und Ultraschall ab.
Rauschen und Mehrwegeffekte bildet die Tiefenkamera nicht ab.

\subsection{Radar und Ultraschall}
\label{sec:starship_radar_us}

Je ein Radar zeigt nach vorne, hinten, links und rechts und sitzt auf 35\,cm Höhe. Radar
erkennt bewegte Objekte wie Radfahrer und Fahrzeuge auch bei Dunkelheit und Regen.
Reichweite und Frequenzbereich orientieren sich an einem 60-GHz-Kurzstreckenradar für
Roboter, dem TI IWR6843AOP \parencite{ti_iwr6843aop}. Dessen Sichtfeld von 120° wurde
bewusst auf 60° $\times$ 20° verengt, damit CARLAs Radar weniger Bodenechos liefert und
die Messpunkte nicht verdünnt. An den Ecken bleiben dadurch Lücken, die die schrägen
Ultraschallsensoren abdecken.

Die acht Ultraschallsensoren bilden einen Ring auf 22\,cm Höhe: drei vorne, drei hinten,
je einer seitlich, wie bei Einparksensoren von Pkw. Die niedrige Lage erfasst auch
niedrige Hindernisse wie Poller oder Tiere. Nachgebildet wird der Ultraschall durch den
CARLA-Hinderniserkennungssensor (Abschnitt~3.2) mit 3\,m Reichweite und einem Suchradius
von 15\,cm; die Reichweite ist an den MaxBotix MB1242 angelehnt \parencite{maxbotix_mb1242}
und wegen der geringen Geschwindigkeit auf 3\,m begrenzt. Die Sensoren sind um 4° nach
oben geneigt; ohne diese Neigung meldeten sie im Test die Fahrbahn als Hindernis. Da der
Sensor nur Treffer meldet, gilt eine Keule als frei, wenn seit 0,3\,s keine Meldung
eingegangen ist.

\subsection{GNSS, IMU und Messraten}
\label{sec:starship_navigation}

Der GNSS-Sensor sitzt oben auf dem Deckel unter der Antenne auf 52\,cm Höhe, die IMU in
Fahrzeugmitte auf 25\,cm. Anders als beim Clevon ist das Rauschen aktiviert: Die
GNSS-Position streut mit einer Standardabweichung von etwa 2\,m, was einem Empfänger ohne
RTK-Korrektur entspricht; die IMU-Werte entsprechen einfachen MEMS-Sensoren. Beide Werte
sind Annahmen ohne Datenblatt des Originals. GNSS liefert 5\,Hz, die IMU ist auf 100\,Hz
eingestellt. Wie in Abschnitt~3.2.1 beschrieben, sind 100\,Hz im asynchronen Betrieb
nicht zugesichert; gemessen wurden je nach Last 28 bis 60\,Hz.

\section{Bordstein-Assistent}
\label{sec:starship_bordstein}

Ein Gehwegroboter muss Bordsteine überqueren. Das reale Fahrzeug klettert mit
Kipphebeln, an denen die Räder sitzen: Die Hebel drehen zunächst frei und werden dann
von einem Motor aktiv gedreht, wobei vor allem die Mittelräder angehoben werden
\parencite{starship_patent_bogie}. Chaos Vehicles kennt keine mechanische Kopplung
zwischen Rädern; jedes Rad ist ein senkrechter Strahl mit eigener Feder.

\subsection{Ausgangsbefund}
\label{sec:starship_bordstein_befund}

Bordsteine in der Referenzkarte Town10HD sind fast überall 15\,cm hoch (Median aus 1150
Messpunkten 14,8\,cm). In der Basismessung an drei festen Stellen scheiterte der Roboter
an allen Bordsteinen beim Hochfahren; er blieb mit der Vorderradmitte etwa 1\,cm vor der
Kante stehen. Herunterfahren gelang, jedoch mit Stoßspitzen von 37 bis 44\,m/s².
Größere Federwege, eine weichere Federung und eine höhere Kollisionsbox änderten die
Stoppposition nicht. Erst mit dreifachem Motormoment stiegen die Vorderräder auf die
Kante; der Roboter kippte dann aber nach hinten über, weil die stark gestauchten
Vorderradfedern die Front weiter anhoben. Eine reine Parameteränderung reichte damit
nicht aus.

\subsection{Umsetzung}
\label{sec:starship_bordstein_umsetzung}

Der Bordstein-Assistent ersetzt den Hebelmotor durch eine geregelte Hubkraft je Rad. Er
ist als C++-Komponente \texttt{UStarshipCurbAssistComponent} im CARLA-Plugin umgesetzt,
an \texttt{BP\_Starship} angehängt und immer aktiv. Die Komponente läuft in jedem
Spieltakt vor der Physikberechnung und benötigt keinen externen Client. Fahrzeugmasse,
Geschwindigkeit und Geometrie bleiben unverändert.

\begin{itemize}
  \item \textbf{Erkennen:} Ein waagerechter Strahl 5\,cm über dem Boden sucht vor jeder
  Fahrzeugseite nach einer Kante. Bei einem Treffer wird die Höhe dahinter gemessen;
  Stufen von 4 bis 25\,cm gelten als Bordstein. Ein Profil im 1-cm-Raster bestimmt, ab
  wo die Oberseite eben ist.
  \item \textbf{Anheben:} Ein PD-Regler hebt das jeweils nächste Rad vor der Kante auf
  eine Sollhöhe knapp über der Oberkante. Die Kraft wirkt am einzelnen Rad, links und
  rechts getrennt, und ist auf 150\,N je Rad begrenzt.
  \item \textbf{Reihenfolge:} Zuerst wird die Vorderachse angehoben, danach Mittel- und
  Hinterachse. Die Sollhöhen liegen auf einer Geraden, sodass das bereits überquerte Rad
  aufsetzt und Vortrieb behält, während das folgende Rad angehoben wird.
  \item \textbf{Schutz:} Ab der Erkennung ist die Geschwindigkeit auf 0,6\,m/s begrenzt;
  über 6° Rollwinkel wird die höhere Seite entlastet, über 40° Nickwinkel die Hubkraft
  der Vorderachse zurückgenommen.
\end{itemize}

Ein Python-Client mit derselben Logik diente als Vorstufe. Er lief asynchron zum Server,
wodurch die Ergebnisse zwischen Durchläufen streuten; die C++-Komponente ist im
Spieltakt reproduzierbar. Eine zusätzliche Anhaltelogik vor der Kante verschlechterte das
Ergebnis auf 0 von 12 Auffahrten und wurde verworfen.

\subsection{Ergebnis und Grenzen}
\label{sec:starship_bordstein_ergebnis}

Tabelle~\ref{tab:starship_bordstein} fasst die Messung aus vier Durchläufen an den drei
Stellen der Basismessung zusammen. Geradeausfahrt, Höchstgeschwindigkeit und
Bremsweg blieben unverändert. Die Werte wurden noch mit der Raycast-Radabfrage
ermittelt; eine Messreihe mit der inzwischen verwendeten Spherecast-Abfrage steht aus.
Im Fahrbetrieb quert der Roboter Bordsteine mit Spherecast nach qualitativer
Beobachtung leichter.

\begin{table}[htbp]
  \centering
  \caption{Bordsteinüberquerung mit und ohne Assistent (Town10HD, 15\,cm, senkrechte
  Anfahrt mit Vollgas, Raycast-Radabfrage); eigene Messung vom 20.10.2026.}
  \label{tab:starship_bordstein}
  \begin{tabular}{lll}
    \hline
    Prüfung & ohne Assistent & mit Assistent \\
    \hline
    Bordstein hoch & 0 von 3 & 10 von 12 (2,7 bis 6,0\,s inkl. 1,2\,m Anlauf) \\
    Bordstein runter & 2 von 2 & 8 von 8 \\
    Stoßspitze beim Herunterfahren & 37 bis 44\,m/s² & 30 bis 46\,m/s² \\
    Nickwinkel beim Hochfahren & -- & 27 bis 33° \\
    Rollwinkel & -- & 1 bis 6° \\
    \hline
  \end{tabular}
\end{table}

Ein starrer Aufbau mit 38\,cm Radstand muss sich an einer 15\,cm hohen Kante
geometrisch um mindestens 28° neigen; das Original vermeidet dies durch seine
Gelenkhebel. Etwa jede sechste Auffahrt bleibt hängen, wenn das Mittelrad die Kante zu
tief erreicht; kurzes Zurücksetzen und erneutes Anfahren behebt dies. Hindernisse bis
etwa 25\,cm Höhe werden als Bordstein erklettert. Das Herunterfahren erfolgt ohne
Unterstützung mit harter Landung. Der Assistent ist eine Nachbildung der Wirkung, kein
Nachbau der Starship-Mechanik.
